\section{Nilpotent cotangent cones and transverse modifications}
\label{sec:geometry}

The detection arguments require two geometric facts.  The nilpotent cone
must have at most half the dimension of the cotangent bundle on a smooth
chart, and a covector outside that cone must admit a Hecke modification
with a nonzero component in the direction of the moving point.  We prove
both facts, including the nondegeneracy of the modification that will be
needed for the slicing argument.

Throughout this section the ground field is algebraically closed.  It is
either of characteristic zero or is $k=\overline{\mathbb F}_q$, with the
four characteristic hypotheses of Theorem~\ref{thm:main}.  In characteristic
zero we choose an invariant form $\kappa$ with the same nondegeneracy
properties.  The stack $\cB$ will be an unlevelled bundle stack, a bundle
stack with ordinary or enhanced Borel level at one point, or, in
characteristic zero, the simultaneous torus quotient
$[(\cA_1^+\times\cA_2^+)/T]$ described in Section~\ref{sec:foundations}.
An actual Hecke modification is always made on one curve, away from its
marked point.

\subsection{Cotangent vectors and the dimension bound}

Let $\mathcal P$ be a $G$-bundle, with reduction to a subgroup $H$ at $x$,
where $H=B$ or $R_u(B)$, and put $\mathfrak h=\Lie(H)$.
The vector bundle governing its deformations is
the elementary modification
\[
 \mathcal E_H
 =\ker\bigl(\ad(\mathcal P)\longrightarrow
                  \ad(\mathcal P)_x/\mathfrak h\bigr).
\]
Its first cohomology is the space of infinitesimal deformations and its
zeroth cohomology is the infinitesimal automorphism space.  Serre duality
therefore identifies a degree-zero cotangent vector with a section
\begin{equation}\label{eq:geometry-cotangent}
 \phi\in H^0\bigl(X,\ad(\mathcal P)\otimes\omega_X(x)\bigr),
 \qquad \operatorname{Res}_x\phi\in\mathfrak h^\perp.
\end{equation}
Here and below a cotangent vector on a smooth stack means an element of
$H^0$ of the fiber of its cotangent complex.  Thus these descriptions
retain the infinitesimal automorphisms of the bundle stack; they do not
replace the stack by a coarse moduli space.

The $T$-weight decomposition and invariance of $\kappa$ show that opposite
root spaces pair perfectly and that $\kappa$ is nondegenerate on
$\mathfrak t$.  Consequently
\[
 \mathfrak b^\perp=\Lie R_u(B),\qquad
 (\Lie R_u(B))^\perp=\mathfrak b.
\]
Formula~\eqref{eq:geometry-cotangent} is thus the strongly parabolic
residue condition at ordinary level and the Borel residue condition at
enhanced level.  For the simultaneous torus quotient, its cotangent
vectors are pairs of enhanced-level Higgs fields annihilating the
infinitesimal diagonal torus action.  Under Serre duality the functional
on a change of enhancement is pairing with the toral component of the
residue, with the sign fixed by the action convention on that branch.

Write $\Lambda\subset T^*\cB$ for the cone of such fields that are
nilpotent over the function field of each curve.  It is closed: after
choosing homogeneous generators of the invariant polynomials, the
condition is the vanishing of their associated global sections.  The
identification of this zero fiber with generic nilpotence is valid under
Hypothesis~\ref{hyp:chevalley}.  Indeed an element can be conjugated into a Borel Lie
algebra, as also verified in the proof of
Lemma~\ref{lem:central-cocharacter} below; contraction of its nilradical
part leaves its toral part.  Restriction of invariants to the torus and
the finite Weyl group quotient then show that all positive-degree
invariants vanish exactly when that toral part is zero.  A finite-group
quotient separates its geometric orbits in every characteristic; no
averaging over the Weyl group is involved.

The residue of a generically nilpotent field is nilpotent.  For example,
in a parameter $t$ at $x$, write $\phi=b(t)\,dt/t$.  Every invariant
polynomial vanishes on $b(t)$ and hence on $b(0)$.  At enhanced level
$b(0)\in\mathfrak b$, so its toral projection vanishes.  It follows that
the enhanced-level cone is the smooth pullback of the ordinary-level
cone.  The same is true for the simultaneous torus quotient over
$\cA_1\times\cA_2$: both toral residues are already zero on its
nilpotent cone.  These assertions also describe the cones after
pullback to scheme charts.

\begin{proposition}\label{prop:cone-dimension}
Let $\cB$ be one of the bundle stacks just specified, and let
$f:D\to\cB$ be a smooth chart with $D$ a smooth finite-type scheme of
pure dimension $d$.  For the smooth cotangent pullback
$\Lambda_D=f^\circ\Lambda\subset T^*D$, one has
\begin{equation}\label{eq:geometry-cone-dimension}
                         \dim\Lambda_D\leq d.
\end{equation}
In characteristic zero this cone is isotropic, in the sense that the
symplectic form restricts to zero on the smooth locus of every reduced
subvariety of $\Lambda_D$.
\end{proposition}

\begin{proof}
We give the ordinary-level argument; the unlevelled version omits the
condition at $x$.  This is the parabolic form of the isotropy argument
for the global nilpotent cone \cite[Lemma~7]{GinzburgNilpotent}; compare the
separable-lifting criterion and its bundle-stack application in
\cite[Appendix~D.1.5 and D.1.8]{AGKRRV}.

Let $Z$ be an irreducible component of $\Lambda_D$ with its reduced
structure, and put $F=k(Z)$.  Its generic point specifies a family
$(\mathcal P,\beta,\phi)$ on $X_F$.  By
\cite[Theorem~2]{DrinfeldSimpson}, after a finite separable extension
$F'/F$ the bundle is trivial over the generic point of $X_{F'}$.
Hypothesis~\ref{hyp:nilpotent}, applied to the field $F'(X)$, gives a parabolic subgroup
whose nilradical contains the value of $\phi$ there.  In characteristic
zero the same assertion is the rational parabolic consequence of
Jacobson--Morozov.  Choose the standard parabolic $P$ of this type.
For the split group $G$, the scheme of parabolics of that type is
$G/P$; the rational parabolic therefore gives, in the generic
trivialization, a section of $\mathcal P/P$ over $F'(X)$.
This is the required generic $P$-reduction and does not require a
choice of rational conjugating element.  Properness of $G/P$ extends it to a reduction
$\mathcal Q$ on the whole smooth curve $X_{F'}$.  Since its nilradical
bundle is a subbundle of $\ad(\mathcal P)$, the generic containment
extends to
\[
 \phi\in H^0\bigl(X_{F'},
          \mathcal Q\mathbin{\times}^{P}\Lie R_u(P)
                     \otimes\omega_{X_{F'}}(x)\bigr).
\]

Consider the stack $\mathcal Y$ of a $P$-bundle together with a Borel
flag at $x$ in the relative position of $(\mathcal Q_x,\beta)$ just
obtained.  This stack is smooth.  Indeed $\Bun_P$ is smooth by the
vanishing of the second cohomology of its deformation bundle on a
curve, and the relative-position locus over it is the associated
bundle with fiber a $P$-orbit in $G/B$.  That orbit is smooth: its
stabilizer is an intersection of a parabolic and a Borel, with the
usual smooth torus and root-group description.

The pullback of $\phi$ to a cotangent vector on $\mathcal Y$ is zero.
To check this without suppressing the level condition, deformations
on $\mathcal Y$ are governed by the locally free sheaf $\mathcal E$
of sections of $\ad(\mathcal Q)$ whose value at $x$ lies in the
intersection with the Borel Lie algebra.  Pairing $\phi$ with
$\mathcal E$ has no pole at $x$, by the original residue condition.
It is zero generically because the nilradical of a parabolic is the
annihilator of its Lie algebra.  Hence the induced section of
$\mathcal E^*\otimes\omega$ is zero everywhere.  Serre duality proves
the asserted vanishing on deformation spaces.

Shrink to a smooth dense open of $Z$.  The finite separable extension
$F'/F$ spreads to a finite \emph{\'etale} cover of a smaller such open,
and the reduction spreads to a map from that cover to $\mathcal Y$.
The canonical one-form of $T^*D$ pulls back to zero there: it evaluates
the given covector on the induced deformation of the bundle, and this
deformation factors through $\mathcal Y$.  Since the cover is separable,
the one-form already vanishes on the smooth dense open of $Z$.
Its exterior differential, the symplectic form, also vanishes there.
An isotropic tangent space in $T^*D$ has dimension at most $d$, proving
\eqref{eq:geometry-cone-dimension}.  The separability in this argument
is essential; a purely inseparable cover would not detect vanishing
of differential forms.

Products and smooth pullbacks preserve this dimension bound: on
cotangent bundles a smooth pullback adds exactly the relative dimension
to the dimension of the cone.  The description of the enhanced and
quotient cones above therefore proves the remaining cases.  In
characteristic zero the same argument can start with any reduced
irreducible subvariety of the cone, rather than only an irreducible
component.  It gives the stated isotropy, which is likewise preserved
by products and smooth pullbacks.
\end{proof}

\subsection{A central cocharacter detecting a nonnilpotent value}

The modification below must be defined by an integral cocharacter.
Choosing merely an element of a toral Lie algebra would lose
information in positive characteristic.  The next lemma supplies an
actual cocharacter and the formal normal form needed for the residue
calculation.

\begin{lemma}\label{lem:central-cocharacter}
Over an algebraically closed field of characteristic zero, or over
$k=\overline{\mathbb F}_q$ under the four characteristic hypotheses,
let $a(t)\in\mathfrak g[[t]]$ with $a(0)$ not nilpotent.
After changing the formal $G$-frame, there are a Levi subgroup $M$ and
a cocharacter $\lambda:\mathbb G_m\to Z(M)$ such that
\begin{equation}\label{eq:geometry-levi-normal-form}
 \begin{gathered}
 a(t)\in\mathfrak m[[t]],\qquad
 \ad(a(0)):\mathfrak g/\mathfrak m\xrightarrow{\sim}
                                      \mathfrak g/\mathfrak m,\\
                         \kappa(a(0),d\lambda)\ne0.
 \end{gathered}
\end{equation}
Here $d\lambda$ is the image of $1\in\Lie\mathbb G_m$ under the
differential of $\lambda$.
\end{lemma}

\begin{proof}
We construct the Levi from the semisimple part of $a(0)$, choose a
central cocharacter that pairs nontrivially with that part, and then
remove the remaining off-Levi coefficients of the series.  The first
two steps require some care in positive characteristic.

\emph{The semisimple part and its centralizer.}
Embed $G$ as a
closed subgroup of some $\operatorname{GL}(W)$.  In positive
characteristic $\mathfrak g$ is closed under the restricted $p$-power,
which is matrix $p$-th power in this representation.  All eigenvalues
of $a(0)$ belong to a finite subfield of $k$.  For a sufficiently large
power $p^r$ that fixes these eigenvalues and kills every nilpotent
Jordan block, matrix power $a(0)^{p^r}$ is its semisimple part $s$.
Consequently both $s$ and $n=a(0)-s$ belong to $\mathfrak g$ and
$[s,n]=0$.  In characteristic zero this is the usual Jordan
decomposition in an algebraic Lie algebra.

To identify $s$ and $n$ intrinsically, fix a maximal torus $T$, a simple system
$\Delta$, and compatible root vectors $e_\alpha,e_{-\alpha}$ with
$[e_\alpha,e_{-\alpha}]=d\alpha^\vee$.  Invariance gives
\begin{equation}\label{eq:geometry-root-pairing}
 \kappa(d\alpha^\vee,H)=c_\alpha\,d\alpha(H),
 \qquad c_\alpha=\kappa(e_\alpha,e_{-\alpha})\ne0
 \quad(H\in\mathfrak t).
\end{equation}
The nonzero constant follows from the perfect pairing of opposite
$T$-weight spaces.  Hypothesis~\ref{hyp:form}, applied to the Levi $T$, makes
$\kappa|_{\mathfrak t}$ nondegenerate.  Since $G$ is simply connected,
the simple coroots form a basis of $X_*(T)$; their differentials are
therefore a basis of $\mathfrak t$.  Equation~\eqref{eq:geometry-root-pairing}
shows that the simple-root differentials are linearly independent.
The same equation, applied to any root, shows that its differential
is nonzero.

We claim that matrix semisimplicity agrees with the toral definition
in the theorem, and that matrix nilpotence agrees with the orbit-closure
definition.  The proper incidence morphism
\[
             G\mathbin{\times}^{B}\mathfrak b\longrightarrow\mathfrak g
\]
is surjective.  Its differential is surjective at a toral element
outside the finitely many root-differential hyperplanes, so its closed
image contains a dense open.  Inside $\mathfrak b$, write an element
as $H+u$.  Successive conjugations by root groups, in increasing root
height, remove every root coefficient with $d\alpha(H)\ne0$:
the coefficient in question changes by a nonzero multiple of the
root-group parameter, while changes to other coefficients occur in
higher heights.  The remaining term $u_0$ commutes with $H$ and is a
nilpotent matrix.  Thus $H+u_0$ is matrix semisimple only if $u_0=0$.
Likewise a matrix nilpotent in $\mathfrak b$ has zero toral projection,
since a faithful representation is upper triangular on $B$ and its
toral differential is injective.  Contraction by a strictly dominant
cocharacter sends the nilradical to zero.  This proves that matrix
nilpotence is equivalent to the orbit-closure definition; the reverse
implication also follows directly from the closed matrix nilpotent
cone.  The same argument applies inside any Levi.

We may now conjugate $s$ into $\mathfrak t$.  Since $a(0)$ is not
nilpotent, $s\ne0$.  Hypothesis~\ref{hyp:centralizer} makes its scheme-theoretic
centralizer a Levi $M=Z_G(s)$.  We have
$s\in\mathfrak z(\mathfrak m)$ and $n\in\mathfrak m$.
Conjugate $n$ within $M$ into a Borel nilradical of $\mathfrak m$,
using the preceding argument.  Torus weights then show that $n$ is
orthogonal to $\mathfrak z(\mathfrak m)$; this conclusion is unchanged
by the conjugation.

\emph{Choosing an integral central cocharacter.}
To detect $s$ by a cocharacter, we need the differentials of central
cocharacters to span $\mathfrak z(\mathfrak m)$.  We verify this before
using the nondegeneracy of $\kappa$ on the center.
Choose the simple system $\Delta$ so that $M$ corresponds to a subset
$\Delta_M\subset\Delta$.
Its Lie center is
\[
 \mathfrak z(\mathfrak m)
       =\bigcap_{\alpha\in\Delta_M}\ker(d\alpha)\subset\mathfrak t.
\]
There are no additional central root vectors, since each root
differential is nonzero.  The integral map
\[
 X_*(T)\longrightarrow\mathbb Z^{\Delta_M},
 \qquad \nu\longmapsto
                 (\langle\alpha,\nu\rangle)_{\alpha\in\Delta_M}
\]
has full row rank after reduction modulo $p$ in positive characteristic,
by the linear independence of the simple-root differentials proved above.
Its Smith normal form therefore has no nonunit elementary divisor
divisible by $p$.
It follows that the kernel of its differential is spanned over $k$
by the differentials of its integral kernel.  The latter kernel is
exactly $X_*(Z(M))$.  Thus the Lie center is spanned by differentials
of actual central cocharacters.  In characteristic zero the same
conclusion follows by tensoring the integral kernel with $k$.

By Hypothesis~\ref{hyp:form} the restriction of $\kappa$ to
this center is nondegenerate, and by the integral-kernel calculation
its central cocharacter differentials span it.  Some
$\lambda\in X_*(Z(M))$ therefore satisfies
\[
             \kappa(a(0),d\lambda)
                    =\kappa(s,d\lambda)\ne0.
\]
On $\mathfrak g/\mathfrak m$ the operator $\ad(s)$ is invertible and
$\ad(n)$ is nilpotent and commutes with it.  Their sum is invertible,
giving the middle assertion of
\eqref{eq:geometry-levi-normal-form}.

\emph{Putting the formal series in the Levi.}
It remains to put the whole series in $\mathfrak m[[t]]$.  Decompose
$\mathfrak g$ into the zero and nonzero eigenspaces of $\ad(s)$,
so that the first summand is $\mathfrak m$.  Suppose inductively that
the non-$\mathfrak m$ coefficients below order $t^r$ have been removed.
Conjugation by a group element congruent to $1+t^rY$ modulo $t^{r+1}$
changes the offending coefficient by $[Y,a(0)]$.  The invertibility
just proved supplies $Y$ that removes it.  Such a group element exists
by smoothness of $G$.  The successive changes converge in $G(k[[t]])$,
proving the first assertion without changing the constant term.
\end{proof}

\subsection{The two tangent maps of a Hecke modification}

The central-cocharacter modification below follows the microlocal method of
\cite[Sections~20.5 and~20.7--20.8]{AGKRRV} and the transverse formulation
in \cite[Proposition~4.2]{OpenAIFullSupport2026}. We include both tangent
calculations, since the slicing argument uses them on opposite fibers.

Let $U$ denote the allowed moving-point curve, or the disjoint union of
the two allowed curves in the quotient case.  An exact-relative-position
Hecke correspondence $H$ has maps
\[
 p:H\longrightarrow\cB,\quad o:H\longrightarrow\cB,\quad
 \operatorname{leg}:H\longrightarrow U,
 \qquad q=(o,\operatorname{leg}),\quad
 \widetilde p=(p,\operatorname{leg}).
\]
Thus $p$ remembers the input bundle and $o$ the output bundle.
For relative position $\lambda^+$, both $q$ and $\widetilde p$ are
representable smooth of relative dimension
$m=\langle2\rho,\lambda^+\rangle$.  This is the usual smooth
affine-Grassmannian orbit description in disc frames.  Its Schubert
closure is a bound proper over either bundle together with the leg;
the inverse modification gives the description from the other side.
These facts hold in families of smooth legs, and a level condition at
a disjoint point is transported unchanged.

\begin{proposition}\label{prop:transverse-hecke}
Let $A\in T^*_{b}\cB$ be a covector outside $\Lambda$.  There exist a
point $y\in U$, an exact-relative-position correspondence $H$, a
modification $h\in H$ with $p(h)=b$ and leg $y$, an output covector
$A'\in T^*_{o(h)}\cB$, and a nonzero $\xi\in T^*_yU$ such that
\begin{equation}\label{eq:geometry-covector-identity}
                         p^*A=q^*(A',\xi)\quad\text{at }h.
\end{equation}
Put
\[
 H_{\mathrm{in}}=\widetilde p^{-1}(b,y),\qquad
 H_{\mathrm{out}}=q^{-1}(o(h),y),
\]
where the fibers include the fixed-side bundle identifications.  Both
are smooth of dimension $m$.  The following two sections have
nondegenerate zeros at $h$:
\begin{enumerate}
\item on $H_{\mathrm{in}}$, the restriction of $p^*A$ to
      $T_{H/(\cB\times U),q}|_{H_{\mathrm{in}}}$;
\item on $H_{\mathrm{out}}$, the restriction of $o^*A'$ to
      $T_{H/(\cB\times U),\widetilde p}|_{H_{\mathrm{out}}}$.
\end{enumerate}
Each is a section of the dual of the indicated rank-$m$ bundle.
In particular, both zeros are isolated.
\end{proposition}

\begin{proof}
Represent $A$ by its Higgs field or pair of fields.  On a curve where
it is not generically nilpotent, choose an unmarked point $y$ at which
the value is not nilpotent.  In a parameter $t$ centered at $y$ and a
formal frame, write this field as $a(t)\,dt$.
Lemma~\ref{lem:central-cocharacter} supplies $M$ and $\lambda$ satisfying
\eqref{eq:geometry-levi-normal-form}.  Make the modification by
$\lambda(t)$, with the convention that its adjoint output lattice is
\[
 L'=\operatorname{Ad}(\lambda(t))L,
 \qquad L=\mathfrak g[[t]],\qquad J=L\cap L'.
\]
Because $\lambda$ centralizes $M$, the field $a(t)\,dt$ is regular in
both lattices.  It therefore transports to a global output covector
$A'$ with the same residue conditions at every marked point.

We record explicitly the tangent spaces and their pairing.  Gauge
transformations on the input formal disc move $L'$, whereas those on the output
move $L$.  Quotienting in each case by the common stabilizer gives
\begin{equation}\label{eq:geometry-two-tangents}
 T_hH_{\mathrm{in}}=L/J,
 \qquad T_hH_{\mathrm{out}}=L'/J.
\end{equation}
Both $L$ and $L'$ are their own annihilators for the pairing
$\operatorname{Res}_{t=0}\kappa(-,-)\,dt$.  More explicitly, if
$\mathfrak g=\bigoplus_{j\in\mathbb Z}\mathfrak g_j$ is the
\emph{integer} weight decomposition for $\lambda$, then
\begin{equation}\label{eq:geometry-truncated-lattices}
 \begin{aligned}
 L/J&=\bigoplus_{j>0}
           \mathfrak g_j\otimes k[[t]]/(t^j),\\
 L'/J&=\bigoplus_{j<0}
           \mathfrak g_j\otimes t^jk[[t]]/k[[t]].
 \end{aligned}
\end{equation}
Opposite weights pair perfectly under $\kappa$, and the coefficient
of $t^i$ pairs with that of $t^{-1-i}$.  Thus
\begin{equation}\label{eq:geometry-residue-perfect}
 (L/J)\times(L'/J)\longrightarrow k,
 \qquad (D,C)\longmapsto\operatorname{Res}\kappa(D,C)\,dt
\end{equation}
is a perfect pairing of $m$-dimensional vector spaces.

To compute the covector identity, first hold the output and leg fixed.
A tangent vector represented by $C\in L'$ changes the input bundle
by that principal part.  Serre duality evaluates $p^*A$ on it as
$\operatorname{Res}\kappa(a(t),C)\,dt$, which vanishes because
$a(t)\in L'$.  The degree-zero smooth cotangent exact sequence for
$q$ implies that $p^*A$ is the pullback of a covector on output bundle
times leg.  This argument is valid for stack covectors: the pullback
in that exact sequence is injective, and its image consists precisely
of covectors annihilating the relative tangent space.

The bundle component is $A'$.  It can be tested on principal parts
at points away from $y$ and the markings, where the input and output
are identified.  Such tests span the deformation space: sufficiently
large pole divisors supported at those points kill the first
cohomology of the deformation bundle, and the principal-parts exact
sequence then surjects onto that first cohomology.  For the moving-point
component, hold the output fixed and replace $t$ by $t-z$ in the
modification.  The logarithmic derivative of $\lambda(t-z)$ at $z=0$
is $-d\lambda/t$.  Consequently, up to the common sign determined by
the gluing convention,
\begin{equation}\label{eq:geometry-leg-covector}
                 \xi=\pm\kappa(a(0),d\lambda)\,dz\ne0.
\end{equation}
These computations prove \eqref{eq:geometry-covector-identity}.
In the simultaneous quotient case they may first be made with both
enhancements retained and then descended along the smooth torus
quotient; the modified curve is disjoint from all gluing data.

It remains to prove the two nondegeneracy assertions.  The perfect
pairing above pairs tangents to the two fixed-side fibers; its
$\ad(a(t))$-twist will be the derivative of each section.
Move in $H_{\mathrm{in}}$ in direction $D\in L/J$ and transport a
test vector $C\in L'/J$ with the moving lattice.  Differentiating its
pairing with the fixed input covector gives
\begin{equation}\label{eq:geometry-derivative-pairing}
 \operatorname{Res}\kappa(a(t),[D,C])\,dt
       =\operatorname{Res}\kappa([a(t),D],C)\,dt,
\end{equation}
up to the same harmless choice of sign.  The expression is independent
of lifts of $D$ and $C$: $\ad(a(t))$ preserves $L$, $L'$ and $J$.
It preserves every $\lambda$-weight space because
$a(t)\in\mathfrak m[[t]]$ and $\lambda$ is central in $M$.
For $j\ne0$, $\mathfrak g_j$ lies in the noncentralizer summand
$\mathfrak g/\mathfrak m$; hence the constant term of this operator
is invertible on $\mathfrak g_j$.  It is therefore invertible on
each of the truncated modules in
\eqref{eq:geometry-truncated-lattices}.  The perfect residue pairing
\eqref{eq:geometry-residue-perfect} makes
\eqref{eq:geometry-derivative-pairing} perfect as well.  This is exactly
the derivative of the first section in the statement at its zero.

For the section on $H_{\mathrm{out}}$, keep the output fixed, move
$L$ using $C\in L'/J$, and transport $D\in L/J$.  The resulting derivative is
the transpose of \eqref{eq:geometry-derivative-pairing}, up to sign.
It too is perfect.  A section of a rank-$m$ vector bundle on a smooth
$m$-dimensional fiber whose derivative at a zero is invertible has
an isolated reduced zero there, proving both assertions.  All
truncation lengths in this proof are integer cocharacter weights;
none is inverted in $k$.
\end{proof}

The two fibers will have different roles in the detection argument.
The nondegenerate zero on $H_{\mathrm{out}}$ isolates a contribution
to a proper Hecke pushforward.  The zero on $H_{\mathrm{in}}$ supplies
the coordinates transverse to that contribution and yields the exact
split quadratic equation of Section~\ref{sec:detection}.
